
\begin{bt}
Cho tam giác \(ABC\) có \(AB=7\), \(AC=9\), \(\widehat{A}=60^\circ\). Tính độ dài cạnh \(BC\).
\shortans{}
\loigiai{
Áp dụng định lý cosin:
\[
BC^2=AB^2+AC^2-2\cdot AB\cdot AC\cos A
=7^2+9^2-2\cdot 7\cdot 9\cdot \cos 60^\circ
=49+81-126\cdot \frac12=67.
\]
Suy ra
\[
BC=\sqrt{67}.
\]
}
\end{bt}

\begin{bt}
Cho tam giác \(ABC\) có \(AB=10\), \(AC=6\), \(BC=8\). Tính \(\widehat{A}\).
\shortans{}
\loigiai{
Đặt \(a=BC=8,\; b=CA=6,\; c=AB=10\). Theo định lý cosin:
\[
a^2=b^2+c^2-2bc\cos A
\Rightarrow 64=36+100-2\cdot 6\cdot 10\cos A.
\]
\[
120\cos A=72\Rightarrow \cos A=\frac{3}{5}.
\]
Vậy
\[
\widehat{A}=\arccos\!\left(\frac35\right)\approx 53{,}13^\circ.
\]
}
\end{bt}

\begin{bt}
Cho tam giác \(ABC\) có \(BC=12\), \(\widehat{A}=30^\circ\), \(\widehat{B}=45^\circ\). Tính cạnh \(CA\).
\shortans{}
\loigiai{
Đặt \(a=BC=12,\; b=CA\). Theo định lý sin:
\[
\frac{a}{\sin A}=\frac{b}{\sin B}
\Rightarrow b=a\cdot \frac{\sin B}{\sin A}
=12\cdot \frac{\sin 45^\circ}{\sin 30^\circ}
=12\cdot \frac{\frac{\sqrt2}{2}}{\frac12}
=12\sqrt2.
\]
Vậy \(CA=12\sqrt2\).
}
\end{bt}

\begin{bt}
Cho tam giác \(ABC\) có \(AB=5\), \(AC=7\), \(BC=9\).
\begin{enumerate}
\item Tính \(\widehat{A}\).
\item Tính diện tích tam giác \(ABC\).
\end{enumerate}
\shortans{}
\loigiai{
Đặt \(a=BC=9,\; b=CA=7,\; c=AB=5\).

Theo định lý cosin:
\[
\cos A=\frac{b^2+c^2-a^2}{2bc}
=\frac{49+25-81}{2\cdot 7\cdot 5}
=-\frac{1}{10}.
\]
Suy ra \(\widehat{A}=\arccos(-0{,}1)\approx 95{,}74^\circ\).

Diện tích:
\[
S=\frac12 bc\sin A=\frac12\cdot 7\cdot 5\cdot \sin A.
\]
Mà
\[
\sin A=\sqrt{1-\cos^2A}=\sqrt{1-\frac1{100}}=\frac{3\sqrt{11}}{10}.
\]
Nên
\[
S=\frac{35}{2}\cdot \frac{3\sqrt{11}}{10}
=\frac{21\sqrt{11}}{4}.
\]
}
\end{bt}

\begin{bt}
Cho tam giác \(ABC\) có \(AB=8\), \(AC=10\), \(\widehat{A}=120^\circ\).
\begin{enumerate}
\item Tính \(BC\).
\item Tính bán kính đường tròn ngoại tiếp \(R\).
\end{enumerate}
\shortans{}
\loigiai{
Đặt \(a=BC,\; b=CA=10,\; c=AB=8\).

Theo định lý cosin:
\[
a^2=b^2+c^2-2bc\cos A
=10^2+8^2-2\cdot 10\cdot 8\cos 120^\circ
=100+64-160\cdot\left(-\frac12\right)=244.
\]
\[
a=2\sqrt{61}.
\]

Theo định lý sin mở rộng:
\[
\frac{a}{\sin A}=2R
\Rightarrow
R=\frac{a}{2\sin A}
=\frac{2\sqrt{61}}{2\sin 120^\circ}
=\frac{2\sqrt{61}}{2\cdot \frac{\sqrt3}{2}}
=\frac{2\sqrt{183}}{3}.
\]
}
\end{bt}

\begin{bt}
(Trường hợp mơ hồ của định lý sin) Cho \(a=8,\; b=10,\; \widehat{A}=30^\circ\).
Hỏi có bao nhiêu tam giác thỏa mãn? Tính các góc có thể có của tam giác.
\shortans{}
\loigiai{
Theo định lý sin:
\[
\frac{a}{\sin A}=\frac{b}{\sin B}
\Rightarrow \sin B=\frac{b\sin A}{a}
=\frac{10\cdot \sin 30^\circ}{8}
=\frac{10\cdot \frac12}{8}
=\frac{5}{8}.
\]
Do đó:
\[
B_1=\arcsin\!\left(\frac58\right)\approx 38{,}68^\circ,\qquad
B_2=180^\circ-B_1\approx 141{,}32^\circ.
\]
Kiểm tra:
\[
A+B_1=30^\circ+38{,}68^\circ<180^\circ \Rightarrow \text{nhận}.
\]
\[
A+B_2=30^\circ+141{,}32^\circ<180^\circ \Rightarrow \text{nhận}.
\]
Vậy có \(2\) tam giác.

Các góc còn lại:
\[
C_1=180^\circ-A-B_1\approx 111{,}32^\circ,\quad
C_2=180^\circ-A-B_2\approx 8{,}68^\circ.
\]
}
\end{bt}
